Transport is usually described by the trajectories it produces. This paper describes it by what
survives: a family of admissible subspaces carried along a path, and the gap that decides what can
be resolved at the end of it. Two things follow, and both are theorems rather than slogans.
\emph{Gaps}: a transport invariant is stable to the extent that the gap across the
boundary of the cluster defining it is open. That principle is made quantitative here --- the corpus's own tip-and-tail
law states it, and what is added is how much each level gains. A rank-one projector buys no
stability at all, because in operator norm it moves at the same rate as its eigenvector.
The gain belongs to clusters, whose protection is the external gap. \emph{Totals}: the distance
map of $N$ bodies has a conserved total capacity, $\sum_{i<j}\lVert\nabla d_{ij}\rVert^2=N(N-2)$
for equal unit masses, so the shape only redistributes a total fixed by $N$. On the equal-mass three-body
shape sphere the redistribution is solved exactly: the squared singular values sum to $3$, their
product is $12A^2/(d_{12}^2d_{13}^2d_{23}^2)$, for the map to the distances the degeneracy at the equilateral shape is a linear
isotropic cone of slope $3$ with holonomy $-1$, the rank dropping exactly where the gap is
largest; for the map to the squared distances it is a quadratic fold with holonomy $+1$, and the
two statements are both true of their own operator. Every claim carries a grade, and the finite ones are machine-checked: 29 kernel-checked theorems, 79 gate
assertions, six deliberately false statements that the checker must reject.
